\documentclass{ctexart}

% 常用宏包
\usepackage{amsmath,amssymb}
\usepackage{bm}
\usepackage{siunitx}
\usepackage{graphicx}
\usepackage{booktabs}
\usepackage{float}
\usepackage{caption}
\usepackage{subcaption}
\usepackage{geometry}
\usepackage{hyperref}

% 页面设置
\geometry{
    a4paper,
    left=2.5cm,
    right=2.5cm,
    top=2.5cm,
    bottom=2.5cm
}

\begin{document}

\begin{center}
    {\LARGE 《星系天文学》第2章作业}

    \vspace{1em}

    姓名：何飞杭 \qquad 学号：202631101066
\end{center}

\noindent\textbf{1.	简单地画出一张赫罗图，注明横坐标和纵坐标，并标出恒星的分布特征。}

\vspace{24.25em}

\vspace{0.5em}

\noindent\textbf{2.	一对目视双星的两颗子星：一颗视星等为1\(^\mathrm{m}\)，另一颗视星等为2\(^\mathrm{m}\)，问这个双星系统的星等是多少？}

\vspace{0.25em}
\[m = -2.5 \log_{10}\frac{f_1 + f_2}{f_0} = -2.5 \log_{10}(10^{-0.4 m_1} + 10^{-0.4 m_2}) \approx 0.638\]
\vspace{0.5em}

\noindent\textbf{3.	已测得南门二(半人马座\(\alpha\)星)和心宿二(天蝎座\(\alpha\)星)的目视星等分别为-0.01\(^\mathrm{m}\)和1.0\(^\mathrm{m}\)，并测量得到它们的距离分别为1.34pc和170pc，求它们的光度之比。}

\vspace{0.25em}
\[M_1 = m_1 - 5 \log_{10}(\frac{d_1}{10}) \approx 4.355 \quad M_2 = m_2 - 5 \log_{10}(\frac{d_2}{10}) \approx -5.15\]
\[L_1/L_2 = 10^{0.4(M_2 - M_1)} \approx 10^{0.4(-5.15-4.355)} \approx 10^{0.4(-9.505)} \approx 10^{-3.802} \approx 1.58 \times 10^{-4}\]
\vspace{0.5em}

\noindent\textbf{4.	(1)关于星系的哈勃分类，请查找资料解释什么是早型星系？什么是晚型星系？}

\vspace{0.25em}
早型星系（early-type galaxies）主要是椭圆星系 E 和透镜状星系 S0，通常以球状/核球结构为主、冷气体少、恒星形成活动弱、恒星总体较老；晚型星系（late-type galaxies）主要是旋涡和不规则星系，通常具有明显的盘结构、较丰富的冷气体和尘埃以及更活跃的恒星形成。
\vspace{0.5em}

\noindent\textbf{(2)早型星系是比晚型星系更年轻的星系。请判断并解释理由。}

\vspace{0.25em}
这里的 early/late 是形态序列上的“早/晚”，类似于“序列的前端/后端”，不是形成时间上的早/晚。此外，早型星系通常以老年恒星为主，恒星形成活动较弱；晚型星系含有更多冷气体和尘埃，恒星形成活动更活跃，因此通常包含更多年轻恒星。

\end{document}